\begin{abcliste}
\abc
\begin{align}
 \kappa &= \kaF \\
   &= \frac{\sqrt{2\times\ncme\times\dEX\times\nce}}{\nchbar} \\
   &= \ka \approx \result{\kanmP{0}{2}}
\end{align}
(With $h$ instead of $\hbar$ the result would be $2\pi$ times too small.)

\abc The ratio of the currents is
\begin{align}
 k &= \frac{e^{-2\kappa(d-\Delta d)}}{e^{-2\kappa d}} = \fkF \\
   &= e^{2\times\ka\times\dd} = \fk \approx \result{\fkP{0}{2}}
\end{align}
A tenth of a nanometre, less than the size of an atom, changes the current about tenfold. The microscope keeps the current constant and records how the tip must move: a map of the atoms.

\abc Because the current falls off so steeply (exponentially) with the distance, $\result{\text{almost all of it flows through the one atom of the tip closest to the surface}}$. The tip never touches the surface: the electrons tunnel across the gap. Binnig and Rohrer received the Nobel prize for this microscope in 1986, developed at IBM Zurich.
\end{abcliste}
